% !TEX root =./ECE444_Practice_main.tex
%
% L10 practice set (Patch, Slot, and Horn Antennas) -- sizing a rectangular
% microstrip patch from the transmission-line design set, the substrate
% bandwidth trade, Babinet complementarity for slots, aperture gain and the
% optimum-horn compromise, and choosing among the three for an application.

\fullwidth{\textbf{Learning Objective 2.3: } I can describe the radiation mechanism, pattern, and typical
use cases for patch, slot, and horn antennas.
\\[4pt]\normalfont\small Show your work. A \textbf{Documentation} line at the end
is required (write \textbf{None} if you did not collaborate). Take
$c = 3\times10^{8}\ \text{m/s}$ and $\eta_0 = 377\ \Omega$.}

\begin{questions}

%%%%%%%%%%%%%%%%%%%%%%%%%%%%%%%%%%%%%%%%%%%%%%%%%%%%%%%%%%%%%%%%%%%%%%%%%%%%
\question \textbf{Sizing a patch.} You are designing a rectangular microstrip patch to resonate
at $f_r = 3.5\ \ghz$ on Rogers RO4003C: $\varepsilon_r = 3.55$, substrate thickness
$h = 0.813\ \text{mm}$. The transmission-line design set is
\eq{
W &= \frac{c}{2f_r}\sqrt{\frac{2}{\varepsilon_r+1}}
\qquad
\varepsilon_{\mathrm{eff}} = \frac{\varepsilon_r+1}{2} + \frac{\varepsilon_r-1}{2}
\lp 1 + \frac{12h}{W} \rp^{-1/2} \\[4pt]
\frac{\Delta L}{h} &= 0.412\ 
\frac{(\varepsilon_{\mathrm{eff}}+0.3)(W/h+0.264)}{(\varepsilon_{\mathrm{eff}}-0.258)(W/h+0.8)}
\qquad
L = \frac{c}{2 f_r \sqrt{\varepsilon_{\mathrm{eff}}}} - 2\Delta L
}

\begin{parts}

	\begin{minipage}[t]{0.58\linewidth}
		\part Find the patch width $W$ in millimetres.
		\begin{solution}[1.0in]
		\eq{
		\frac{c}{2f_r} &= \frac{3\times10^{8}}{2(3.5\times10^{9})} = 42.86\ \text{mm} \\
		W &= 42.86\sqrt{\frac{2}{4.55}} = 42.86(0.6630) = 28.4\ \text{mm}
		}
		\end{solution}
	\end{minipage}\hfill%
	\begin{minipage}[t]{0.37\linewidth}
		\raggedleft \vspace{12pt}
		$W$ = \ansbox[1.3in]{$28.4\ \text{mm}$}
	\end{minipage}

	\begin{minipage}[t]{0.58\linewidth}
		\part Find the effective permittivity $\varepsilon_{\mathrm{eff}}$.
		\begin{solution}[1.0in]
		\eq{
		\frac{12h}{W} &= \frac{12(0.813)}{28.41} = 0.3433 \\
		\varepsilon_{\mathrm{eff}} &= 2.275 + 1.275(1.3433)^{-1/2} \\
		&= 2.275 + 1.275(0.8628) = 3.375
		}
		\end{solution}
	\end{minipage}\hfill%
	\begin{minipage}[t]{0.37\linewidth}
		\raggedleft \vspace{12pt}
		$\varepsilon_{\mathrm{eff}}$ = \ansbox[1.3in]{$3.375$}
	\end{minipage}

	\begin{minipage}[t]{0.58\linewidth}
		\part Find the edge extension $\Delta L$ and the physical length $L$.
		\begin{solution}[1.4in]
		With $W/h = 28.41/0.813 = 34.95$,
		\eq{
		\frac{\Delta L}{h} &= 0.412\ \frac{(3.675)(35.21)}{(3.117)(35.75)}
		= 0.412(1.1614) = 0.4785 \\
		\Delta L &= 0.4785(0.813) = 0.389\ \text{mm} \\
		L &= \frac{42.86}{\sqrt{3.375}} - 2(0.389) = 23.33 - 0.78 = 22.55\ \text{mm}
		}
		\end{solution}
	\end{minipage}\hfill%
	\begin{minipage}[t]{0.37\linewidth}
		\raggedleft \vspace{12pt}
		$\Delta L$ = \ansbox[1.2in]{$0.389\ \text{mm}$}
		\\[10pt]
		$L$ = \ansbox[1.2in]{$22.6\ \text{mm}$}
	\end{minipage}

	\part The half-wavelength in the dielectric is $\lambda_d/2 = 23.3\ \text{mm}$, yet you etched
	$L = 22.6\ \text{mm}$. Explain in one or two sentences what physical effect forces the metal
	to be cut short, and state which way the resonant frequency moves if you etch the
	full $23.3\ \text{mm}$ instead.
		\begin{solution}[1.1in]
		The fields fringe past each open edge, so the resonator is electrically longer than the
		copper by $\Delta L$ at each end. Resonance is set by the electrical length, so the physical
		metal must be shortened by $2\Delta L$ to keep the total at $\lambda_d/2$. Etching the full
		$23.3\ \text{mm}$ makes the antenna electrically too long, and the resonance moves
		\emph{down} in frequency --- here by roughly $3\%$, which is several times the patch's own
		bandwidth, so the antenna would be badly mismatched at $3.5\ \ghz$.
		\end{solution}

\end{parts}

\newpage
%%%%%%%%%%%%%%%%%%%%%%%%%%%%%%%%%%%%%%%%%%%%%%%%%%%%%%%%%%%%%%%%%%%%%%%%%%%%
\question \textbf{The substrate trade.} A patch is a high-$Q$ cavity, and its impedance bandwidth
for VSWR $\le 2$ follows the closed form
\eq{
\text{BW} \approx 3.77\ \frac{\varepsilon_r-1}{\varepsilon_r^{2}}\ \frac{h}{\lambda_0}\ \frac{W}{L}
}
Continue with the $3.5\ \ghz$ design of Question 1 ($W = 28.41\ \text{mm}$, $L = 22.55\ \text{mm}$,
$h = 0.813\ \text{mm}$, $\varepsilon_r = 3.55$).

\begin{parts}

	\begin{minipage}[t]{0.58\linewidth}
		\part Estimate the fractional bandwidth of the RO4003C design, as a percentage.
		\begin{solution}[1.1in]
		$\lambda_0 = c/f_r = 85.7\ \text{mm}$, so $h/\lambda_0 = 0.00949$.
		\eq{
		\text{BW} &\approx 3.77\lp\frac{2.55}{12.60}\rp(0.00949)\lp\frac{28.41}{22.55}\rp \\
		&= 3.77(0.2023)(0.00949)(1.260) = 0.0091
		}
		About $0.9\%$ --- roughly $32\ \text{MHz}$ of usable band.
		\end{solution}
	\end{minipage}\hfill%
	\begin{minipage}[t]{0.37\linewidth}
		\raggedleft \vspace{12pt}
		BW = \ansbox[1.3in]{$0.91\%$}
	\end{minipage}

	\begin{minipage}[t]{0.58\linewidth}
		\part Rebuild the same $3.5\ \ghz$ patch on $\varepsilon_r = 10.2$ at the same $h$. Its design
		works out to $W = 18.11\ \text{mm}$, $L = 13.35\ \text{mm}$. Find the new bandwidth, and the
		factor by which the patch \emph{area} shrank.
		\begin{solution}[1.3in]
		\eq{
		\text{BW} &\approx 3.77\lp\frac{9.2}{104.0}\rp(0.00949)\lp\frac{18.11}{13.35}\rp \\
		&= 3.77(0.0884)(0.00949)(1.356) = 0.0043
		}
		Area: $28.41\times22.55 = 641\ \text{mm}^2$ against $18.11\times13.35 = 242\ \text{mm}^2$,
		a factor of $2.6$ smaller. Bandwidth fell by a factor of $2.1$.
		\end{solution}
	\end{minipage}\hfill%
	\begin{minipage}[t]{0.37\linewidth}
		\raggedleft \vspace{12pt}
		BW = \ansbox[1.2in]{$0.43\%$}
		\\[10pt]
		area $\div$ \ansbox[1.2in]{$2.6$}
	\end{minipage}

	\part Read the bandwidth formula: which single substrate parameter would you change to
	\emph{buy} bandwidth without changing $\varepsilon_r$, and name one penalty that comes with it.
		\begin{solution}[1.1in]
		Increase the thickness $h$: bandwidth is directly proportional to $h/\lambda_0$. The penalties
		are surface waves (power trapped in the substrate instead of radiated, so efficiency and
		pattern quality drop), a larger probe or feed inductance that is harder to match, and a
		thicker, heavier board. Above about $h = 0.05\lambda_0$ the closed forms above stop being
		trustworthy at all.
		\end{solution}

	\part You are given $8\ \text{mm} \times 8\ \text{mm}$ of board area for a $3.5\ \ghz$ element
	and told the link needs $2\%$ bandwidth. State whether a single patch on any of these
	substrates can do the job, and what you would tell the systems lead.
		\begin{solution}[1.2in]
		No. Even the low-$\varepsilon_r$ option is far larger than $8\ \text{mm}$, and the only
		substrate that fits the footprint ($\varepsilon_r = 10.2$, at $18\times13\ \text{mm}$) is still
		too big \emph{and} delivers $0.43\%$, not $2\%$. Shrinking with permittivity and widening the
		bandwidth pull in opposite directions on this element. Tell the lead the requirement needs a
		different answer: a much thicker substrate, a stacked or aperture-coupled patch, or a
		different element entirely --- not a tweak to this one.
		\end{solution}

\end{parts}

\newpage
%%%%%%%%%%%%%%%%%%%%%%%%%%%%%%%%%%%%%%%%%%%%%%%%%%%%%%%%%%%%%%%%%%%%%%%%%%%%
\question \textbf{Babinet and the slot.} A narrow half-wave slot is cut in a large, thin
conducting sheet and fed across its middle.

\begin{parts}

	\part State Babinet's complementarity relation between a slot and its complementary dipole,
	and say in one sentence what ``complementary'' means physically.
		\begin{solution}[1.1in]
		\eq{ Z_{\text{slot}}\ Z_{\text{dipole}} = \frac{\eta_0^{2}}{4} }
		The complementary structure is the one you get by exchanging metal for air and air for
		metal: the slot is exactly the hole left behind if the dipole were punched out of the sheet.
		\end{solution}

	\begin{minipage}[t]{0.58\linewidth}
		\part The complementary dipole is trimmed to resonance, where it presents $73\ \Omega$ real.
		Find the slot's input impedance.
		\begin{solution}[1.1in]
		\eq{
		\frac{\eta_0^{2}}{4} &= \frac{(377)^{2}}{4} = 3.553\times10^{4}\ \Omega^{2} \\
		Z_{\text{slot}} &= \frac{3.553\times10^{4}}{73} = 487\ \Omega
		}
		This is the canonical ``about $485\ \Omega$'' resonant-slot value.
		\end{solution}
	\end{minipage}\hfill%
	\begin{minipage}[t]{0.37\linewidth}
		\raggedleft \vspace{12pt}
		$Z_{\text{slot}}$ = \ansbox[1.3in]{$487\ \Omega$}
	\end{minipage}

	\begin{minipage}[t]{0.58\linewidth}
		\part Now take the untrimmed $\lambda/2$ dipole, $Z = 73 + j42.5\ \Omega$. Find the slot
		impedance and comment on the sign of its reactance.
		\begin{solution}[1.4in]
		\eq{
		Z_{\text{slot}} &= \frac{3.553\times10^{4}}{73+j42.5}
		= \frac{3.553\times10^{4}(73-j42.5)}{73^{2}+42.5^{2}} \\
		&= \frac{3.553\times10^{4}}{7135}(73-j42.5) = 4.980(73-j42.5) \\
		&= 364 - j212\ \Omega
		}
		The inductive dipole becomes a \emph{capacitive} slot: inverting a complex impedance flips
		the sign of the reactance. Both therefore go resonant at the same length.
		\end{solution}
	\end{minipage}\hfill%
	\begin{minipage}[t]{0.37\linewidth}
		\raggedleft \vspace{12pt}
		$Z_{\text{slot}}$ = \ansbox[1.5in]{$364 - j212\ \Omega$}
	\end{minipage}

	\part A horizontal slot is cut in the vertical fin of an aircraft. State the polarization of
	the radiated field and justify it in one sentence, then name one reason a designer chooses a
	slot here rather than a monopole.
		\begin{solution}[1.2in]
		The field is \emph{vertically} polarized. The slot's \textbf{E} field runs \emph{across} the
		cut rather than along it (fields swap under Babinet), so a horizontal slot radiates like a
		vertical dipole. The slot is chosen because it is flush: no protrusion, no drag, nothing to
		shear off or ice up, and it can be cut into structure that is already there.
		\end{solution}

\end{parts}

\newpage
%%%%%%%%%%%%%%%%%%%%%%%%%%%%%%%%%%%%%%%%%%%%%%%%%%%%%%%%%%%%%%%%%%%%%%%%%%%%
\question \textbf{An X-band horn.} A pyramidal standard-gain horn has a rectangular aperture
$16.5\ \text{cm} \times 12.2\ \text{cm}$ and is used at $9.4\ \ghz$. Aperture antennas obey
$G = \eta_{ap}\,4\pi A/\lambda^{2}$; take the optimum-horn value
$\eta_{ap} = 0.5$ unless told otherwise.

\begin{parts}

	\begin{minipage}[t]{0.58\linewidth}
		\part Find the gain, in dBi.
		\begin{solution}[1.3in]
		\eq{
		\lambda &= \frac{3\times10^{8}}{9.4\times10^{9}} = 3.19\ \text{cm} \\
		A &= 0.165(0.122) = 0.0201\ \text{m}^{2} \\
		\frac{4\pi A}{\lambda^{2}} &= \frac{4\pi(0.0201)}{(0.0319)^{2}} = 248 \\
		G &= 0.5(248) = 124 = 10\ \mylog{124} = 20.9\dbi
		}
		\end{solution}
	\end{minipage}\hfill%
	\begin{minipage}[t]{0.37\linewidth}
		\raggedleft \vspace{12pt}
		$G$ = \ansbox[1.3in]{$20.9\dbi$}
	\end{minipage}

	\begin{minipage}[t]{0.58\linewidth}
		\part How far away must a receiving antenna sit for this horn to be in its far field?
		\begin{solution}[1.1in]
		The largest aperture dimension is the diagonal,
		\eq{
		D &= \sqrt{16.5^{2}+12.2^{2}} = 20.5\ \text{cm} \\
		r &\ge \frac{2D^{2}}{\lambda} = \frac{2(0.205)^{2}}{0.0319} = 2.6\ \m
		}
		\end{solution}
	\end{minipage}\hfill%
	\begin{minipage}[t]{0.37\linewidth}
		\raggedleft \vspace{12pt}
		$r$ = \ansbox[1.3in]{$2.6\ \m$}
	\end{minipage}

	\begin{minipage}[t]{0.58\linewidth}
		\part A shorter, more sharply flared horn of the same aperture has enough edge phase error
		to drop $\eta_{ap}$ to $0.35$. What gain does it deliver, and how much was lost?
		\begin{solution}[1.1in]
		\eq{
		G &= 0.35(248) = 86.9 = 19.4\dbi \\
		\Delta G &= 10\ \mylog{\frac{0.5}{0.35}} = 1.6\db
		}
		Same physical area, $1.6$\db{} less gain --- all of it spent on phase error.
		\end{solution}
	\end{minipage}\hfill%
	\begin{minipage}[t]{0.37\linewidth}
		\raggedleft \vspace{12pt}
		$G$ = \ansbox[1.2in]{$19.4\dbi$}
		\\[10pt]
		loss = \ansbox[1.2in]{$1.6\db$}
	\end{minipage}

	\part Explain in two or three sentences why enlarging a horn's aperture at fixed flare length
	eventually stops increasing its gain, and what the phrase ``optimum horn'' names.
		\begin{solution}[1.2in]
		The wave leaving the flare is spherical while the aperture is flat, so the aperture edge is
		farther from the virtual apex than the center and its phase lags. Widening the aperture at
		fixed length raises $4\pi A/\lambda^{2}$ but steepens that quadratic phase error, so
		$\eta_{ap}$ falls; past a point the efficiency loses faster than the area gains and the
		curve turns over. The \emph{optimum horn} is that peak --- the shortest horn for a given
		aperture whose edge phase error is still tolerable (about $\lambda/4$ in the E-plane and
		$3\lambda/8$ in the H-plane), which is where $\eta_{ap} \approx 0.5$ comes from.
		\end{solution}

\end{parts}

\newpage
%%%%%%%%%%%%%%%%%%%%%%%%%%%%%%%%%%%%%%%%%%%%%%%%%%%%%%%%%%%%%%%%%%%%%%%%%%%%
\question \textbf{Mechanism and selection.} Answer in words. Full sentences.

\begin{parts}

	\part A patch is mostly a flat sheet of copper over a ground plane, and a sheet of copper is
	a terrible radiator. Explain what actually radiates, and why the two contributions add rather
	than cancel in the broadside direction.
		\begin{solution}[1.4in]
		The cavity between the patch and the ground plane holds a half-wave standing wave, so the
		field points from patch to ground at one open edge and from ground to patch at the other. At
		each open edge the field fringes out past the conductor. The vertical parts of the two
		fringing fields are opposite and cancel; the horizontal parts point the same way and add.
		Those two fringing edges behave as two slots of width $W$, spaced $L \approx \lambda_d/2$
		apart, driven in phase --- so their contributions arrive with equal path length and add
		straight up, giving a broadside beam.
		\end{solution}

	\part The two radiating slots sit only about a third of a free-space wavelength apart. Say
	what that implies for the E-plane beamwidth, and why the patch's directivity therefore lands
	near $6$\dbi{} rather than up with the horns.
		\begin{solution}[1.2in]
		A two-element array that short cannot form a narrow beam: the array factor barely varies
		across the visible half space, so the E-plane pattern stays very broad. The whole radiating
		structure is only a fraction of a wavelength across, and directivity tracks electrical size,
		so a patch radiates a fat hemispherical beam worth roughly $5$ to $8$\dbi. Horns get their
		$20$-plus dBi by being many wavelengths across.
		\end{solution}

	\part Match each application to the best of patch, slot, or horn, with a one-line
	justification: \textbf{(i)} the gain reference in an anechoic chamber; \textbf{(ii)} a
	telemetry antenna on a supersonic missile skin; \textbf{(iii)} the element of a printed
	phased array for a $10\ \ghz$ radar demonstrator.
		\begin{solution}[1.4in]
		\textbf{(i)} Horn --- built to the optimum design and calibrated, so its gain is known to a
		few tenths of a dB, which is exactly what a reference must be.
		\textbf{(ii)} Slot --- it is flush with the skin, so it survives the aerodynamic and thermal
		environment; cavity-backing makes it radiate outward only.
		\textbf{(iii)} Patch --- planar, light, cheap, and reproducible, so hundreds of identical
		elements and their feed network etch in the same step.
		\end{solution}

	\part Your team wants the phased-array demonstrator of part (c) to cover a $10\%$ band. State
	the conflict this creates with the patch element and one concrete way to resolve it.
		\begin{solution}[1.2in]
		A thin-substrate patch delivers a few percent at best, so $10\%$ is roughly an order of
		magnitude beyond a plain element --- and thinning the requirement is not an option if the
		waveform needs the band. Concrete fixes: go to a thick, low-$\varepsilon_r$ substrate;
		stack a parasitic patch above the driven one; or feed through an aperture-coupled slot,
		which adds a second resonance and typically buys $10$ to $20\%$. Each costs board
		thickness, layers, or both --- the bandwidth has to be paid for somewhere.
		\end{solution}

\end{parts}

\vspace{1.5em}
\noindent\textbf{Documentation:} \rule[-2pt]{0.6\linewidth}{0.4pt}

\end{questions}
